\beginsong{Clameurs du monde}[by={François Tardif, Johan Tardif, Maïwenn Tardif}]
\beginverse
\[Bm] Malgré les peurs, les tempêtes,\[G] les appels de la plan\[D]ète,
Un cri\[A] persiste.\[Bm]
Sous les cendres des gu\[G]erres, les droits du monde résis\[D]tent,
Nos c{\oe}u\[A]rs grandissen\[Bm]t.
Que nos clameurs franch\[G]issent les murs et les fro\[D]ntières,
Que les voix des \[A]peuples s'unissen\[Bm]t.
Clameurs de ceux qu'on ou\[G]blie, de ceux qui es\[D]pèrent,
Et pour les démuni\[A]s, rendons justice.\[Bm]
\endverse
\beginverse*
Ensembl\[G]e, on répon\[D]d à l'appel\[A] du mond\[Bm]e,
Ensemb\[G]le, on répon\[D]d à l'appe\[A]l.
\endverse*
\beginchorus
\[Bm]Écoute, \[G]les voix r\[D]ugir du fond\[A] de nos c{\oe}urs !
\[Bm]Écoute, \[G]les voix r\[D]ugir oh, \[A]oh ! 
Faisons monter nos \[Bm]clameurs, \[G]clameurs, \[D]clameurs oh, \[A]oh !
Faisons monter nos \[Bm]clameurs, \[G]clameurs, \[D]clameurs oh, \[A]oh !
\endchorus
\beginverse
\[Bm] Malgré les bouleversements\[G], les changements de la ter\[D]re,
La jo\[A]ie apparaît\[Bm].
Sous le regard silen\[G]cieux, toutes les nations s'affr\[D]ontent,
La pa\[A]ix renaît\[Bm].
Que nos clameurs dé\[G]voilent nos espoirs et nos \[D]rêves,
Que toutes nos \[A]angoisses disparaiss\[Bm]ent.
Clameurs des c{\oe}urs bri\[G]sés, des âmes déchirées\[D],
Et pour l'aven\[A]ir, une promess\[Bm]e.
\endverse
\beginverse
[Pont - version clip]
\[Bm] Écoutez, le v\[G]ent murmure des voix,\[D]
Des milliers de rêves qui s’él\[A]èvent en chœur,
Ils portent l’es\[Bm]poir, l’amour, la j\[G]oie,
Et font tre\[D]mbler les cœurs, ils brisent\[A] la peur.
Car il est \[Bm]temps, ô frères, ô sœurs,\[G] De faire monter les cla\[D]meurs du monde,
Pour que ja\[A]mais, ô jamais la pa\[Bm]ix ne soit qu'un rêve,
Mais une vé\[G]rité qui gronde et in\[D]onde, inonde, i\[A]nonde
\endverse
\endsong
